\section{Implementation blueprint}
\label{app:implementation}

The listings record the numerical core.  They omit only ordinary exception
handling and the choice of quadrature library.

\subsection{Stable grid quantities}

The endpoint gauges should be evaluated from log-odds, not from rounded ranks:
\begin{equation}
\ell_-(\Lambda(z))=1+\log(1+e^{-z}),
\qquad
\ell_+(\Lambda(z))=1+\log(1+e^z).
\label{eq:stablegauges}
\end{equation}
Both logarithms use a softplus routine.  This remains accurate when
$\Lambda(z)$ has rounded to zero or one.

\begin{lstlisting}[caption={Build one generalized LQD slice.}]
def build(theta, alpha, grid):
    # Endpoint chart: log lambda-, right-tail coordinate, body modes.
    lam_m, lam_p, a = decode_endpoint_chart(theta, alpha)
    L, R = raw_chord(lam_m, lam_p, a)

    u = expit(grid.z)
    ell_m = 1 + softplus(-grid.z)
    ell_p = 1 + softplus( grid.z)
    rho = expit(grid.z) * expit(-grid.z)

    g = (1-u)*L + u*R + a @ grid.Pmat
    log_dx = g - alpha.minus*log(ell_m) \
               - alpha.plus *log(ell_p)
    check_exponent_budget(log_dx)
    dx = exp(log_dx)

    if alpha.plus == 0 and lam_p >= 1 - wall_eps:
        raise InadmissibleForward
    if dx[-1] > 1 - tail_eps:
        raise RemoteSaddleBeyondGrid

    xbar = cumulative_simpson(dx, anchor=grid.mid)
    check_exponent_budget(xbar)

    I = integrate(exp(xbar)*rho)
    I += left_tail_mass(xbar[0], theta, alpha, grid)
    I += right_tail_mass(xbar[-1], theta, alpha, grid)
    m = -log(I)
    x = m + xbar

    dG = -exp(x)*rho
    G = reverse_integrate(-dG)
    G += right_tail_mass(x[-1], theta, alpha, grid)
    return Slice(x=x, dx=dx, G=G, dG=dG,
                 m=m, lam=(lam_m, lam_p),
                 alpha=alpha, theta=theta)
\end{lstlisting}

The two tail-mass functions use the continuation
\eqref{eq:rightcontinuation} and its left counterpart.  For
$\alpha=0$ they use \eqref{eq:tailcorr}.  For $\alpha>0$ they integrate in
log space after the change $z=Z_{\max}+t$.  All tail corrections and their
parameter derivatives enter both the primal build and the sensitivity pass.

The median anchor fixes $\bar x(0)=0$.  Simpson quadrature is fourth order on
the uniform grid.  Exact nodal derivatives make the cubic Hermite
representation fourth order off grid without a global spline solve.

\subsection{Call evaluation}

\begin{lstlisting}[caption={Evaluate one normalized call.}]
def call_price(slc, k, grid):
    if k > slc.x[-1]:
        z = right_tail_root(slc, k, grid)  # analytic tail continuation
        share = right_ledger_tail(slc, z, grid)
        cash = exp(k - softplus(z))
        return share - cash

    if k < slc.x[0]:
        put = left_put_tail(slc, k, grid)
        return 1 - exp(k) + put

    z = bracketed_hermite_root(slc.x, slc.dx, k)
    share = hermite_eval(slc.G, slc.dG, z)
    cash = exp(k - softplus(z))
    return share - cash
\end{lstlisting}

The right branch uses the same tail law as the martingale normalizer.  A
separate volatility extrapolation would generally break the density and
calendar identities.

\subsection{Sensitivity pass}

\begin{lstlisting}[caption={All fixed-tail parameter sensitivities.}]
def sensitivities(slc, grid):
    # Rows are derivatives of g in the chosen coefficient chart.
    B = basis_columns_in_chart(slc.theta, grid.u)
    dxbar = cumulative_simpson(slc.dx * B,
                               anchor=grid.mid, axis=1)
    weight = exp(slc.x) * grid.rho
    dm = -integrate(weight * dxbar, axis=1)
    add_tail_derivatives_to_shift(dm, slc, grid)
    dx = dm[:, None] + dxbar

    integrand = weight[None, :] * dx
    dG = reverse_integrate(integrand, axis=1)
    add_tail_derivatives_to_ledger(dG, slc, grid)
    dG_slope = -integrand
    dvs = -2 * integrate(dx * grid.rho, axis=1)
    return Sensitivities(dG=dG, dG_slope=dG_slope, dvs=dvs)

def price_jacobian_row(slc, sens, k):
    z = strike_root(slc, k)
    return hermite_eval(sens.dG, sens.dG_slope, z)
\end{lstlisting}

The martingale derivative is
$\partial m=-\int e^x\partial\bar x\,\rho\,\dd z$, including both tail
pieces.  Envelope cancellation then turns the ledger sensitivity at the
strike root directly into the call sensitivity.

\subsection{Full-line calendar constraints}

During the joint maturity fit, maintain a small active set of ranks for every
adjacent expiry pair.

\begin{lstlisting}[caption={Exchange algorithm for global calendar order.}]
def fit_surface(nodes):
    active = initial_calendar_ranks(nodes)
    while True:
        surface = constrained_joint_solve(nodes, active)
        worst = full_line_calendar_minima(surface)
        if worst.value >= -calendar_tol:
            return surface
        active[worst.pair].add(worst.rank)

def full_line_calendar_minima(surface):
    out = []
    for near, far in adjacent_pairs(surface):
        roots = all_piecewise_cubic_roots(far.x - near.x)
        roots += analytic_tail_roots(far, near)
        points = [-inf] + roots + [inf]
        gaps = [0.0] \
             + [far.ledger(z) - near.ledger(z) for z in roots] \
             + [0.0]
        out.append(min_with_location(gaps, points))
    return min(out)
\end{lstlisting}

If the difference of the two cubic transports vanishes on an entire segment,
the ledger gap is constant there by \eqref{eq:calgapderivative}; either
endpoint represents the segment.  Tail root isolation uses
\eqref{eq:rightcontinuation}.  The final certificate is rerun on the
publication grid and with tighter tail quadrature than the optimizer uses.

\subsection{Failure tests}

The implementation treats the following states separately.
\begin{enumerate}
\item \emph{Rank saturation:} use \eqref{eq:stablegauges}, two-sided logistic
factors, and log-space cash legs.
\item \emph{Interior overflow:} inspect the full arrays of $g$, $\log x'$,
and $\bar x$ before exponentiation; admissible endpoints do not bound body
modes.
\item \emph{Exponential-tail wall:} for $\alpha_+=0$, keep
$\lambda_+\le1-\varepsilon_{\rm wall}$ in a logistic chart.
\item \emph{Remote light-tail saddle:} for $\alpha_+>0$, apply
\eqref{eq:operationaltailguard} or use a saddle-centered quadrature.
\item \emph{Near-zero option prices:} compare calls and calendars in log space
and obtain plotted remote implied variance from the analytic wing law.
\item \emph{Calendar gaps:} a finite penalty is not an acceptance test; the
quantile-root certificate must pass for every adjacent expiry pair.
\end{enumerate}

\subsection{Reference defaults}

\begin{center}
\small
\begin{tabular}{@{}llll@{}}
\toprule
Quantity & Default & Quantity & Default\\
\midrule
Basis order $N$ & 16 (cap 24) & Order guard & \eqref{eq:orderguard}\\
Grid & $[-40,40]$ & Nodes & 8001 (fit: 2001)\\
Tail exponents & stated scenarios & Exponential wall & $1-10^{-6}$\\
Remote-saddle margin & $10^{-4}$ & Exponent budget & 700\\
Ridge $(\omega_{\rm reg},\nu)$ & $(10^{-6},1)$ & Vega floor $\eta$ & $10^{-4}$\\
Haircut $h$ & 0.005 volatility & Solver & constrained trust region\\
Slice check & 801 strikes & Calendar check & all quantile crossings\\
Optimization tolerance & $10^{-10}$ & Publication rebuild & fourfold check\\
\bottomrule
\end{tabular}
\captionof{table}{Reference numerical defaults.  Tail exponents and their maturity
policy are model inputs and are therefore reported with the fit.}
\label{tab:defaults}
\end{center}

The shortest porting sequence is: reproduce the logistic reference law;
reproduce a normal law with $\alpha_\pm=1/2$; implement the stable cash leg
and both tail continuations; add the analytic sensitivity pass; then add the
one-slice and full-line calendar certificates.
